\section{A Static Slicer for Qiskit Programs}\label{sec:tool}

A slice helps developers only as lines of their program, so we built a tool that slices a Qiskit program
statically with the QDG of Section~\ref{sec:dependence}. Fig.~\ref{fig:tool} shows its steps on the ansatz of
Fig.~\ref{fig:ansatz}.

\paragraph*{Parsing}
The tool parses the program with srcML~\cite{collard2011lightweight, collard2013srcml}, which marks up its
syntax and the position of every element. A table of \numToolGateTable{} Qiskit instruction methods records
which arguments of each are qubits, bits and angles, so that a \texttt{call} element such as
\texttt{qc.cx(qubit, qubit + 1)} becomes a gate on its qubits, as stage~1 of Fig.~\ref{fig:tool} shows.

\paragraph*{Static interpretation}
The tool turns the srcML tree into an operation trace, as Behler et al.\ do for
OpenQASM~\cite{behler2024supporting}. It unrolls loops and records the iteration of each operation, inlines
functions, and follows \texttt{if\_test} and \texttt{c\_if}, Qiskit's constructs for gates conditioned on
measured bits, as stages 2 and 3 of Fig.~\ref{fig:tool} show. Each value carries the lines that computed it,
which gives the classical lines of each operation. Where a
value cannot be resolved statically, the tool over-approximates. An unknown qubit index acts on every qubit
it may denote, an unknown condition runs every branch, and a loop over an unknown range repeats its body
once more than the number of qubits and bits it touches, at most \numToolMaxCopies{} times. Since the needed set never loses a qubit, these over-approximations
only enlarge a slice, and the tool reports each. A circuit passed to code the tool does not model is reported as not traced.

\paragraph*{Slicing}
The slice is the backward pass of Theorem~\ref{thm:sound} on the trace, which skips barriers and terminal
measurements. An operation that may act on several qubits adds all of them to the needed set, and a bit
leaves the needed set only when an operation certainly writes it. The pass returns the operations of the
slice, their lines and iterations, and the classical lines they depend on, in time linear in the trace.

\paragraph*{Forward slices and profiles}
A forward pass over the same trace gives forward slices, dual to the backward ones on our test programs,
and a qubit slice profile combines one pass of each. Over the
\numProfBugsQubits{}, \numProfTextQubits{} and \numProfLintqQubits{} qubits of the Bugs4Q programs, the
textbook and \numProfLintqSample{} sampled LintQ files, the two slices of a qubit differ in length for
\numProfBugsFwdNeBwd{}, \numProfTextFwdNeBwd{} and \numProfLintqFwdNeBwd{} of them.

\paragraph*{Validation}
On \numToolDevPrograms{} runnable programs that we wrote, the static trace equals the circuit Qiskit builds. On \numToolDevSliceEqual{} of them, every slice equals
the circuit-level QDG slice of the study. The remaining program nests conditional blocks, and on
\numToolNestedDraws{} random product inputs the oracle of Section~\ref{sec:design} finds the circuit-level slice
wrong \numToolNestedStudyWrong{} times and the tool's slice never. The tool's line slices equal those of the
study on all \numToolTTwoUnits{} units whose lines the study maps, which include the \numLinesUnits{} units of
Section~\ref{sec:results}. Of the \numToolBugsQiskit{} Qiskit files of Bugs4Q, the \numToolTextbookQiskit{} notebooks of the Qiskit
textbook~\cite{qiskittextbook} and the \numToolLintqQiskit{} programs of LintQ's
corpus~\cite{paltenghi2024lintq}, \numToolBugsCircuit{}, \numToolTextbookCircuit{} and
\numToolLintqCircuit{} yield a circuit, \numToolBugsResolved{} of the Bugs4Q ones without over-approximation.
No Bugs4Q or LintQ program crashes the tool, which takes \numToolBugsMedianMs{}~ms per Bugs4Q program in the median.
